\section{Algoritmet e gjenerimit dhe kampionimit}
\subsection{Gjeneratori dhe gjendja e tij}
Në NumPy rekomandohet objekti \texttt{Generator}. Ai e bën burimin e rastësisë të dukshëm dhe shmang varësinë nga një gjendje globale.
\begin{lstlisting}[language=Python]
rng = np.random.default_rng(2026)
u = rng.random(1_000_000)
print(u.mean(), u.var())  # priten afersisht 1/2 dhe 1/12
\end{lstlisting}

\subsection{Transformimi i anasjelltë}
\begin{lstlisting}[language=Python,caption={Mostra eksponenciale pa përdorur funksionin e gatshëm.}]
def eksponenciale_anasjelltas(rng, lam, size):
    u = rng.random(size)
    return -np.log1p(-u) / lam

x = eksponenciale_anasjelltas(rng, lam=1.5, size=100_000)
print(x.mean(), 1/1.5)
\end{lstlisting}
Funksioni \texttt{log1p(-u)} është më i saktë kur $u$ është i vogël. Kontrollohen mesatarja, varianca dhe CDF-ja empirike.

\subsection{Box--Muller}
\begin{lstlisting}[language=Python,caption={Gjenerimi i çifteve Gauss standarde.}]
def box_muller(rng, n):
    u1 = rng.random(n)
    u2 = rng.random(n)
    r = np.sqrt(-2.0*np.log(u1))
    phi = 2.0*np.pi*u2
    return r*np.cos(phi), r*np.sin(phi)
\end{lstlisting}
Testohet jo vetëm histogrami i secilit komponent, por edhe kovarianca ndërmjet tyre.

\subsection{Kampionimi me refuzim}
\begin{lstlisting}[language=Python,caption={Refuzimi për një dendësi në interval të fundmë.}]
def refuzim(rng, f, a, b, M, n):
    pranuar = []
    prova = 0
    while len(pranuar) < n:
        x = rng.uniform(a, b)
        y = rng.uniform(0.0, M)
        prova += 1
        if y <= f(x):
            pranuar.append(x)
    return np.asarray(pranuar), n/prova
\end{lstlisting}
Këtu $M$ është kufi i sipërm i $f$ në $[a,b]$. Në formulimin me propozim të përgjithshëm, kontrollohet $f(x)\le M g(x)$.

\subsection{Testet minimale}
\begin{enumerate}
\item krahasohen momentet empirike me ato teorike;
\item vizatohet CDF-ja empirike;
\item testohet autokorrelacioni;
\item përsëritet me fara të ndryshme;
\item matet norma e pranimit dhe kostoja për mostër.
\end{enumerate}
Termi \emph{kampionim} përdoret për sampling dhe \emph{mostër} për sample. \emph{Gjenerator pseudorastësor} thekson natyrën deterministe të vargut.

Terminologjia probabilitare shqip mbështetet te \cite{capollari,butka,tasho}; algoritmet e gjenerimit dhe testimit krahasohen me trajtimet standarde \cite{press,robert}.
